\section{Methodology}
This study documents the design and operation of the PaperCircle discovery agent. The system is built on the \texttt{smolagents} framework, which provides the baseline agent primitives (tools, tool-calling agents, and an orchestrator). PaperCircle extends this baseline with a stateful, multi-agent research pipeline that retrieves, ranks, and summarizes academic papers while producing structured artifacts at every step. A central \texttt{PipelineState} object stores the shared state, including the current query, a deduplicated paper list, step logs, and summary statistics. After each agent action, the pipeline updates synchronized outputs (e.g., \texttt{papers.json}, \texttt{links.json}, \texttt{stats.json}, \texttt{summary.json}) and exports auxiliary formats (CSV, BibTeX, Markdown, and a live HTML dashboard).

\subsection{Baseline: smolagents}
The baseline for orchestration is the \texttt{smolagents} library. The pipeline uses a \texttt{CodeAgent} as the orchestrator and multiple \texttt{ToolCallingAgent} instances, each attached to a specific tool. The baseline responsibilities include (i) tool invocation, (ii) multi-step planning via the orchestrator, and (iii) delegation to specialized agents. PaperCircle keeps the baseline behavior intact but adds structured outputs, offline search, and evaluation metrics. The baseline tool interface is preserved: each tool receives explicit parameters and returns a formatted string response, allowing the orchestrator to chain steps while maintaining readability and traceability.

\subsection{System Architecture}
The pipeline is composed of six agents: intent classification, paper search, sorting, analysis, export, and web search. The intent classifier parses natural-language queries into structured constraints (search mode, conferences, year range, max results, and ranking preferences). The paper search agent is the primary retrieval worker; it updates the global state and writes outputs after every search step. The sorting and analysis agents operate on the shared paper list to refine ranking and derive insights. The export agent centralizes output access for downstream workflows, while the web search agent supplements the pipeline with external lookup tools when required. All agents are coordinated by the \texttt{CodeAgent}, which enforces a minimal-step policy for efficiency and uses the intent classifier to decide offline versus online search.

\subsection{State Management and Outputs}
State is maintained in \texttt{PipelineState}. Each step increments a counter, logs action metadata, and regenerates synchronized artifacts. The outputs include: (i) \texttt{papers.json} with full paper metadata and computed scores, (ii) \texttt{links.json} with structured links and PDFs/DOIs, (iii) \texttt{stats.json} with aggregate statistics and a leaderboard, (iv) \texttt{summary.json} with generated insights and key findings, (v) \texttt{retrieval\_metrics.json} when evaluation is enabled, and (vi) human-readable exports (CSV, BibTeX, Markdown) plus a live HTML dashboard. This approach ensures that each agent step is reproducible and auditable.

\subsection{Retrieval}
The pipeline supports both offline and online retrieval. Offline search loads papers from a local JSON corpus and optionally filters by conference and year. It ranks results using BM25 by default, with optional semantic similarity (sentence transformers) or hybrid scoring when available. An optional cross-encoder reranker can refine the top results; when enabled, it reranks a first-stage candidate set. Online search aggregates results from arXiv, Semantic Scholar, OpenAlex, and DBLP via their public APIs. A query intent classifier detects search mode, conference constraints, year ranges, and ranking preferences, and routes the query to the appropriate retrieval pathway. Deduplication is applied across sources by normalizing titles.

\subsection{Ranking and Scoring}
After retrieval, papers are scored along multiple axes: recency, similarity to the query (TF--IDF when available), novelty based on title token frequency, and normalized BM25 scores. The system supports sorting by any single criterion or by a weighted combined score. Relevance scores are computed as a weighted mixture of similarity, recency, citation count, and BM25. Final ranks are assigned after sorting, and the updated ordering is reflected in all exported artifacts. When reranker-based sorting is requested, a cross-encoder replaces the default scoring with direct relevance scores.

\subsection{Analysis and Monitoring}
The pipeline computes aggregate statistics such as source distribution, year distribution, top authors and venues, keyword frequency, and citation summaries. These analytics populate structured summaries and are visualized in an auto-refreshing HTML dashboard. Each agent action is logged with timestamps and paper counts, enabling reproducibility and step-level auditing of the pipeline. The pipeline also maintains a step log that captures the agent name, action, results preview, and parameters used.

\subsection{Evaluation}
When a ground-truth paper title or identifier is provided, the pipeline computes retrieval metrics at each step, including Mean Reciprocal Rank (MRR), Recall@K, Precision@K, and hit rates. These metrics are stored in \texttt{retrieval\_metrics.json} and can be aggregated for benchmarking. A separate benchmarking routine supports parallel evaluation across multiple queries, reporting mean metrics and timing statistics.
